% SOURCE_STATEMENT: "At every perturbative order, is the renormalized planar connected two-point function $G_{\\alpha,\\beta}$ a polynomial with rational coefficients in $\\alpha$, $\\beta$, $A$, $B$, and iterated integrals indexed by rooted trees?"
% Preparation and final self-audit: CONTRIBUTING.md prompts C and D; source/audit artifacts retained in the local physics_statements_2026_10_11 batch.
% !TeX program = lualatex
\documentclass[11pt]{article}
\usepackage[margin=25mm]{geometry}
\usepackage{fontspec}
\setmainfont{Latin Modern Roman}
\usepackage{amsmath,amssymb,mathtools}
\usepackage{unicode-math}
\setmathfont{Latin Modern Math}
\usepackage{xurl,hyperref}
\hypersetup{colorlinks=true,urlcolor=blue}
\setcounter{secnumdepth}{0}
\setlength{\parindent}{0pt}
\setlength{\parskip}{5pt}
\setlength{\emergencystretch}{3em}
\newcommand{\sref}[2]{\hyperlink{source:#1:#2}{[S#2]}}
\newcommand{\cataloguescope}[1]{\paragraph{Recorded target.}#1}
\newcommand{\R}{\mathbb R}
\newcommand{\C}{\mathbb C}
\newcommand{\N}{\mathbb N}
\newcommand{\Z}{\mathbb Z}
\newcommand{\Q}{\mathbb Q}
\newcommand{\F}{\mathbb F}
\newcommand{\eps}{\varepsilon}
\begin{document}
% UnsolvedMath ID 30001336; source catalogue code OWR-4084-010
\setcounter{section}{30001335}
\section{Rooted-Tree Expansions of Renormalized Two-Point Functions}\label{30001336}
{\small UnsolvedMath ID 30001336 · Mathematical Physics}
\subsection{Definitions and mathematical statement}

\paragraph{Shared notation.}$\mathbb N=\{1,2,\ldots\}$; the standard number fields have their usual meanings. All additional conventions are stated locally.


\paragraph{Definitions and conventions.} The source model is the self-dual noncommutative $\phi^4_4$ theory in its planar sector. Its renormalized connected two-point function is expanded formally in the coupling $\lambda$, $G_{\alpha\beta}=\sum_{n\ge0}\lambda^nG^{(n)}_{\alpha\beta}$, with $0\le\alpha,\beta<1$. Write $A=(1-\alpha)/(1-\alpha\beta)$ and $B=(1-\beta)/(1-\alpha\beta)$. Rooted trees index nested integrals arising in the source's integral-equation recursion. The first examples are $I(z)=\int_0^1 z(1-zx)^{-1}\,dx=-\log(1-z)$ and $I_\bullet(z)=\int_0^1 zI(x)(1-zx)^{-1}\,dx$. Polynomiality at every order means that each coefficient $G^{(n)}$ lies in the algebra over $\mathbb Q$ generated by the specified variables and source iterated integrals.

\paragraph{Source scope and qualifications.} The source supplies the first integral examples and checks low orders; it does not give a complete recursive definition for every allowed tree integral in this brief passage. Any all-order proof must fix that source family, including permitted endpoint constants, rather than allow arbitrary integrals that trivialize the assertion. The same conjecture recurs at ID 30001596.

\paragraph{Catalogue mathematical target.}
\paragraph{Catalogue wording retained for traceability.}
At every perturbative order, is the renormalized planar connected two-point function $G_{\alpha,\beta}$ a polynomial with rational coefficients in $\alpha$, $\beta$, $A$, $B$, and iterated integrals indexed by rooted trees?

\paragraph{Saved literature context.} The supplied catalogue's assessment dated 2026-08-21 records: The cited work derives the Schwinger-Dyson equation and verifies the rooted-tree form only through low perturbative orders; the all-order assertion remains a conjecture. The assessment is retained as dated provenance; it is not a new comprehensive literature or open-status certification.
\subsection{Short English statement}
Prove the rooted-tree iterated-integral form of every perturbative coefficient of the planar two-point function.
\subsection{Sources}
\begin{itemize}
\item[\hypertarget{source:30001336:1}{S1}] ULAM.AI and the UnsolvedMath Contributors, UnsolvedMath, record 30001336 (OWR-4084-010). Catalogue curation under CC BY 4.0; definitions and English exposition added locally.
\url{https://huggingface.co/datasets/ulamai/UnsolvedMath}.
\item[\hypertarget{source:30001336:2}{S2}] Noncommutative Geometry (2009), Oberwolfach Reports, DOI 10.4171/owr/2009/41, pp. 2271--2334; official report source and mathematical context.
\url{https://doi.org/10.4171/owr/2009/41}.
\item[\hypertarget{source:30001336:3}{S3}] Official EMS Press PDF of the same Oberwolfach report; PDF pages 50--51, printed pages 2320--2321.
\url{https://ems.press/content/serial-article-files/46243}.
\item[\hypertarget{source:30001336:4}{S4}] noncommutative field theory preprint; reference in the saved dated catalogue literature review.
\url{https://arxiv.org/abs/0909.1389}.
\item[\hypertarget{source:30001336:5}{S5}] author manuscript; reference in the saved dated catalogue literature review.
\url{https://ivv5hpp.uni-muenster.de/u/raimar/publications/preprints/0909.1389.pdf}.
\end{itemize}
\label{end:30001336}
\end{document}
